% ============================================================
%  chapters/01_introduction.tex
% ============================================================

\section{Introduction}
\label{sec:introduction}

Close to one-third of all nations have experienced some form of internal
conflict in the past 40 years, with roughly twenty percent experiencing a
conflict lasting more than ten years \citep{blattman2009civil}. Moreover,
the poorest nations---those in Sub-Saharan Africa, Asia, Eastern Europe,
and South America---tend to have the highest propensity for war and violent
conflict. These conflicts devastate the lives, health, and living standards
of the people who experience them, even though there has been no conclusive
evidence that in the long run these episodes of violence affect the
development trajectory of a country \citep{davis2002bones,
miguel2011long}.

The purpose of this paper is to measure the effects of violence on
educational attainment among school-age children in Nigeria. I use survey
data from the 2010 Nigeria Education Data Survey (NEDS) and combine this
with violence data spanning 1999--2010 from the Uppsala Conflict Data
Program (UCDP) to examine whether experiencing violence in a child's local
government area (LGA) affects the number of years of education a student
will attain and whether it increases the likelihood that a student will be
enrolled in school.

The history of Nigeria is one marked by violence and unrest since the
country gained independence from Great Britain in 1960. This legacy has
been attributed to the fact that Nigeria's borders, like those of most
former African colonies, were drawn arbitrarily without consideration of the
tribal or ethnic relations of those who would soon call themselves fellow
countrymen. Nigeria's case is particularly complex due to its high degree of
ethnic and tribal diversity: there are more than 300 different ethnic groups
within its borders. Major struggles for control and power arise as tribes
fight for control of natural resource-rich lands, political parties compete
for control of those resources, and religious groups attempt to impose their
ideologies on members of other religious communities.

There is a growing literature that attempts to measure the role violence
plays in educational attainment and economic returns to schooling. In
general, it has been found that violence has a negative effect on
educational attainment and that children exposed to violence earlier in life
attain less schooling than their unexposed peers \citep{leon2012, shemyakina2011}. 
Experiencing violence also appears to have a larger negative effect on girls
than on boys \citep{shemyakina2011}. I review this literature in
Section~\ref{sec:literature}.

The remainder of this paper is organized as follows. Section~\ref{sec:background}
provides a history of education in Nigeria. Section~\ref{sec:literature}
reviews the existing empirical literature on the determinants of school
dropout and the effects of violence on educational attainment.
Section~\ref{sec:data} describes the datasets used in this paper.
Section~\ref{sec:methodology} details the econometric methodology.
Section~\ref{sec:results} presents and discusses the results.
Section~\ref{sec:conclusion} concludes with policy recommendations and
directions for future research.
